\section{Security Analysis and Discussion}
\label{sec:discussion}

This section provides a formal security analysis of \textsc{WasmSentinel}'s isolation mechanisms and discusses its limitations and future directions.

\subsection{Security Analysis}

\textbf{Sandbox Isolation (G2).}
\textsc{WasmSentinel} achieves robust isolation of untrusted detection logic through the WebAssembly (Wasm) sandbox.
Unlike traditional native plugins, Wasm modules execute in a restricted linear memory space, completely decoupled from the host's address space.
Any memory access beyond the allocated bounds triggers an immediate trap, preventing buffer overflow attacks from compromising the monitoring daemon.
Furthermore, we employ the WebAssembly System Interface (WASI) to enforce a capability-based security model.
By default, detection rules are granted \textit{zero} capabilities; they cannot access the filesystem, network, or environment variables.
The only communication channel is the restricted host-call interface for receiving security events and emitting alerts.

\textbf{Resource Exhaustion Protection.}
To mitigate Denial-of-Service (DoS) attacks by malicious or buggy rules, we utilize Wasmtime's fuel-based execution metering.
Each detection module is allocated a fixed amount of ``fuel'' (instruction budget) per event.
If a rule enters an infinite loop or exceeds its budget, the runtime terminates its execution, ensuring the availability of the monitoring system even under adversarial conditions.
Memory usage is similarly capped at 256MB per instance.

\textbf{eBPF Integrity and TCB.}
The integrity of event collection relies on the Linux kernel's eBPF verifier, which ensures that loaded programs are type-safe and terminate.
While historical vulnerabilities in the verifier exist~\cite{ebpf-bugs}, it remains the industry-standard TCB for kernel-level observability.
\textsc{WasmSentinel} minimizes the TCB by keeping the eBPF logic simple (event emission only) and shifting complex detection logic to the Wasm sandbox.

\subsection{Limitations and Evasion}

\textbf{Mimicry and Stealthy Attacks.}
Sophisticated adversaries may attempt to evade sequence-based detection by interleaving malicious syscalls with benign ones (mimicry attacks) or by slowing down the attack to exceed the STIDE window.
Our hybrid architecture partially mitigates this by escalating suspicious (though not definitively anomalous) sequences to the MLP for deeper semantic analysis.
However, a determined adversary with full knowledge of the deployed model's weights could potentially craft adversarial syscall sequences.

\textbf{Detection Scope.}
Currently, \textsc{WasmSentinel} monitors syscall-level activity.
While syscalls are the primary interface for resource access, attacks that operate entirely in user-space memory (e.g., pure computational logic or certain library-level exploits) may remain invisible.
Future extensions could integrate library-call hooking or hardware performance counters to broaden the visibility.

\subsection{Modern Dataset Generalization}
Our evaluation on ADFA-LD and modern CVE-2024+ simulations (Section~\ref{sec:evaluation}) demonstrates that patterns learned from older datasets can generalize to modern container escapes.
This suggests that while specific attack implementations evolve, the fundamental ``syscall signatures'' of privilege escalation and unauthorized resource access remain remarkably consistent over time.

\subsection{Future Work}
We plan to explore federated learning to allow multiple organizations to collaboratively update detection models without sharing sensitive telemetry data.
Additionally, we are investigating the use of Wasm SIMD extensions to further accelerate neural network inference for even more complex models like Transformers.
